\begin{table}[H]
  \centering
  \small
\caption{Description textuelle du cas d’utilisation : «~Créer un utilisateur~»}
  \label{tab:cu-gerer-utilisateurs}
  \PFETableSetup
  \PFETableRowColors
  \PFETableArrayStretch
  \PFETableTabularXBegin
  \begin{tabularx}{\PFETableBlockWidth}{@{}>{\raggedright\arraybackslash}p{3.2cm} >{\raggedright\arraybackslash}X@{}}
    \tableHeaderStyle \tableHeaderCell{Champ} & \tableHeaderCell{Contenu} \\
    \textbf{Titre} & Créer un utilisateur \\
    \textbf{Acteurs} & Utilisateur Corporate \\
    \textbf{Objectif} &
      Permettre à l’utilisateur Corporate de créer des comptes utilisateurs afin de gérer l’accès à la plateforme. \\
    \textbf{Préconditions} &
      1. L’utilisateur Corporate est authentifié.\par
      2. Un token est obtenu pour accéder aux ressources protégées. \\
    \textbf{Postconditions} &
      1. Les informations du nouvel utilisateur sont enregistrées.\par
      2. Le compte utilisateur est disponible dans la liste des utilisateurs. \\
    \textbf{Scénario nominal} &
      1. L’utilisateur Corporate accède au module de gestion des utilisateurs. \par
      2. Le système affiche la liste des utilisateurs existants. \par
      3. L’utilisateur choisit l’action : créer un utilisateur. \par
      4. L’utilisateur saisit les informations du nouvel utilisateur (nom, email, rôle, site associé, mot de passe…). \par
      5. Le système valide les données saisies. \par
      6. Si les données sont valides, le système enregistre la création de l’utilisateur et met à jour la liste des utilisateurs. \\
    \textbf{Scénario alternatif} &
      \textbf{A. Email déjà utilisé :} \par
      A.1 L’utilisateur Corporate saisit une adresse email déjà existante. \par
      A.2 Le système rejette l’opération. \par
      A.3 Un message est affiché : «~Un utilisateur avec cet email existe déjà~». \par
      \textbf{B. Mot de passe non sécurisé :} \par
      B.1 L’utilisateur Corporate saisit un mot de passe ne respectant pas les critères. \par
      B.2 Le système rejette l’opération. \par
      B.3 Un message est affiché : «~Le mot de passe doit contenir au moins 8 caractères, une majuscule, une minuscule et un caractère spécial~». \\
  \end{tabularx}%
  \PFETableTabularXEnd
\end{table}